
\section{Appendix D. Negative Results}

\subsection{D.1 Attempts That Failed}

\subsubsection{D.1.1 Head Pruning Instead of Dim Reduction}

\textbf{Goal.} Prune entire heads instead of reducing head dimension.

\textbf{Result.} Accuracy dropped by {{PRUNE\_DROP}}\% -- more than dim reduction.

\textbf{Conclusion.} Head-dim reduction preserves more information.

\subsubsection{D.1.2 Random Head-Dim Slice}

\textbf{Goal.} Slice random head dimensions instead of leading.

\textbf{Result.} Accuracy dropped by {{RANDOM\_DROP}}\%.

\textbf{Conclusion.} Leading dimensions are more informative.

\subsubsection{D.1.3 No LoRA Recovery}

\textbf{Goal.} Apply head-dim reduction without LoRA.

\textbf{Result.} Accuracy dropped by {{NOLORA\_DROP}}\%.

\textbf{Conclusion.} LoRA is essential for recovery.

\subsubsection{D.1.4 Per-Head LoRA}

\textbf{Goal.} Separate LoRA adapter per head.

\textbf{Result.} No improvement over shared LoRA; more parameters.

\textbf{Conclusion.} Shared LoRA suffices.

\subsubsection{D.1.5 Full Fine-Tuning}

\textbf{Goal.} Train full model after head-dim reduction.

\textbf{Result.} Similar to LoRA but 100$\times$ more compute.

\textbf{Conclusion.} LoRA is more efficient.

\subsection{D.2 Abandoned Variants}

\subsubsection{D.2.1 Learned Head-Dim}

Instead of fixed reduction, learn per-head dimension. Complexity not worth it.

\subsubsection{D.2.2 Adaptive LoRA Rank}

Per-layer rank allocation. Marginal improvement.

\subsubsection{D.2.3 Teacher Distillation}

Distill from larger model. Orthogonal to our method.

\subsection{D.3 Failed Hyperparameters}

\subsubsection{LoRA Rank Too Small}
\begin{table}[H]
\centering
\small
\begin{tabular}{ll}
\toprule
\textbf{$r_L$} & \textbf{DistilBERT} \\
\midrule
1 & {{D\_LR1}} \\
2 & {{D\_LR2}} \\
4 & {{D\_LR4}} \\
\bottomrule
\end{tabular}
\end{table}

\subsubsection{Head-Dim Reduction Too Large}
\begin{table}[H]
\centering
\small
\begin{tabular}{ll}
\toprule
\textbf{$s$} & \textbf{DistilBERT} \\
\midrule
3.0 & {{D\_S30}} \\
4.0 & {{D\_S40}} \\
\bottomrule
\end{tabular}
\end{table}

\subsubsection{Recovery Too Short}
\begin{table}[H]
\centering
\small
\begin{tabular}{ll}
\toprule
\textbf{Steps} & \textbf{DistilBERT} \\
\midrule
5 & {{D\_S5}} \\
10 & {{D\_S10\_STEP}} \\
\bottomrule
\end{tabular}
\end{table}

\subsection{D.4 Fixes That Didn't Help}

\begin{table}[H]
\centering
\small
\begin{tabular}{lll}
\toprule
\textbf{Fix} & \textbf{Problem} & \textbf{Result} \\
\midrule
Per-head LoRA & Quality & No improvement \\
Adaptive rank & Quality & No improvement \\
Random slice & Quality & Worse \\
Weight decay tuning & Quality & No improvement \\
\bottomrule
\end{tabular}
\end{table}

\subsection{D.5 Multi-Seed}

\begin{table}[H]
\centering
\small
\begin{tabular}{ll}
\toprule
\textbf{Seed} & \textbf{DistilBERT} \\
\midrule
42 & 90.25 \\
123 & 90.10 \\
456 & 90.35 \\
\textbf{Mean} & \textbf{90.23} \\
\textbf{Std} & \textbf{0.13} \\
\bottomrule
\end{tabular}
\end{table}

\subsection{D.6 Lessons Learned}

\begin{enumerate}
\item \textbf{Head-dim reduction > head pruning.}
\item \textbf{Leading dimensions are best.}
\item \textbf{LoRA is essential.}
\item \textbf{Rank 16 suffices.}
\item \textbf{Shared LoRA is sufficient.}
\end{enumerate}
\begin{lstlisting}

---
\end{lstlisting}